\begin{frame}{Git, GitHub, GitLab khác nhau thế nào?}
\begin{itemize}
\item Git là công cụ quản lý phiên bản phân tán, chạy trên máy cá nhân.
\item GitHub và GitLab là máy chủ chứa repo, cộng thêm cộng tác và CI/CD.
\item Một project dùng Git ở local, đẩy lên GitHub hoặc GitLab để chia sẻ.
\item Hiểu đúng: Git là nền tảng bắt buộc, GitHub và GitLab là lựa chọn bổ sung.
\end{itemize}
{\small
\begin{tabular}{@{}p{2.0cm}p{3.4cm}p{4.2cm}@{}}
\toprule
Khía cạnh & Git & GitHub, GitLab \\
\midrule
Bản chất & Tool dòng lệnh & Dịch vụ hosting và cộng tác \\
Nơi chạy & Máy cá nhân & Cloud hoặc self-host \\
Lịch sử & Đầy đủ trong .git & Nhân bản cộng thêm PR, MR, CI \\
Ví dụ & \texttt{git log} tại local & Mở PR trên web để review \\
\bottomrule
\end{tabular}}
\end{frame}

\begin{frame}{GitHub và GitLab: 8 tiêu chí cho đồ án}
\begin{itemize}
\item Bảng dưới so sánh 8 điểm ảnh hưởng trực tiếp tới đồ án ASPX.
\item Nhóm môn học ưu tiên: miễn phí, CI dễ học, review nhanh.
\item Cả hai đều hỗ trợ repo private miễn phí cho sinh viên.
\item Chốt nhanh: chọn một nền tảng chính, mirror sang nền tảng còn lại.
\end{itemize}
{\footnotesize
\renewcommand{\arraystretch}{0.9}
\begin{tabular}{@{}p{2.6cm}p{3.3cm}p{3.7cm}@{}}
\toprule
Tiêu chí & GitHub & GitLab \\
\midrule
Giá sinh viên & Free và Education pack & Free và Ultimate giáo dục \\
CI cho ASPX 4.8.1 & Actions, runner Windows & Pipeline, runner Windows \\
Review code & PR và rulesets & MR và approval rules \\
Quản lý việc & Issues, Projects & Issues, Boards, Epics \\
Bảo mật & Dependabot, CodeQL & SAST, secret detection \\
Self-host & Không, chỉ cloud & Có, bản community \\
Tài liệu mẫu & Rất nhiều ví dụ & Đầy đủ, gắn DevOps \\
Phù hợp ASPX & Tốt, Actions dễ học & Tốt, pipeline mạnh \\
\bottomrule
\end{tabular}}
\end{frame}

\begin{frame}[fragile]{Nên chọn nền tảng nào cho đồ án?}
\begin{itemize}
\item Khuyến nghị 1: chọn GitHub làm chính nếu nhóm mới học CI lần đầu.
\item Khuyến nghị 2: chọn GitLab làm chính nếu muốn pipeline nhiều stage và review chặt.
\item Khuyến nghị 3: dùng chính một nơi, mirror sang nơi còn lại để dự phòng.
\item Mọi ví dụ ASPX 4.8.1 đều build bằng runner Windows và \texttt{msbuild}.
\item Thầy cô chấm nhanh nhờ link PR hoặc MR cộng với pipeline xanh.
\end{itemize}
\begin{lstlisting}[language=bash]
# mirror GitHub -> GitLab (backup)
git remote add gitlab <gitlab-url>
git push gitlab main --tags
\end{lstlisting}
\end{frame}

\begin{frame}[fragile]{Quy trình 3 người: main, develop, feature}
\begin{itemize}
\item Luật 1: cấm push thẳng \texttt{main} và \texttt{develop}, mọi đổi phải qua PR hoặc MR.
\item Luật 2: mỗi PR cần một người review và pipeline xanh mới được merge.
\item Luật 3: đặt tên nhánh \texttt{feature/ten-viec}, commit theo Conventional Commits.
\item TV1 giữ \texttt{develop} ổn định, TV2 và TV3 làm trên nhánh feature riêng.
\item Conflict file \texttt{Web.config}: người tạo PR tự merge bản mới nhất rồi sửa.
\end{itemize}
\begin{lstlisting}[language=bash]
# so do nhanh: main <- develop <- feature
main    :  o--------o--------o  (release)
develop :  o---o---o----o---o  (integration)
feature :   \o/      \o/  (1 branch / person)
\end{lstlisting}
\end{frame}

\begin{frame}[fragile]{Review checklist và quản lý secret}
\begin{columns}[t]
\column{0.48\textwidth}
Checklist review:
\begin{itemize}
\item Code build sạch, không warning mới, test chính pass.
\item Không còn marker conflict, đã tự review diff trước khi nhờ review.
\item Cập nhật ghi chú ngắn để thầy cô hiểu đổi gì, vì sao.
\end{itemize}
\column{0.48\textwidth}
Quản lý secret:
\begin{itemize}
\item Không commit connection string thật, mật khẩu, file publish.
\item Dùng biến môi trường và secret của CI cho \texttt{Web.config} khi build.
\item Lỡ lộ secret: thu hồi ngay, xóa lịch sử, báo thầy cô.
\end{itemize}
\end{columns}
\begin{lstlisting}[language=bash]
# remove leaked config from history
git rm --cached Web.config
git commit -m "fix: stop tracking local config"
\end{lstlisting}
\end{frame}

\begin{frame}[fragile]{Bài tập về nhà}
\begin{itemize}
\item Bài 1 (Git local): tạo nhánh \texttt{feature/them-trang}, commit 3 lần, merge vào \texttt{develop}.
\item Bài 2 (PR và CI): mở PR, nhờ bạn review, pipeline build \texttt{msbuild} phải xanh.
\item Bài 3 (Secret): chuyển connection string vào secret của CI, chứng minh lịch sử sạch.
\item Nộp link PR hoặc MR kèm ảnh pipeline xanh trước buổi sau.
\item Chấm điểm: đúng quy trình 40, CI xanh 30, review chất 30.
\end{itemize}
\begin{lstlisting}[language=bash]
# starter for homework
git checkout -b feature/them-trang
git log --oneline -5
git push -u origin feature/them-trang
\end{lstlisting}
\end{frame}